\section{Keterkaitan Buku Ajar dengan RPS Berbasis OBE}

Buku ajar ini dirancang selaras dengan Rencana Pembelajaran Semester (RPS) mata kuliah Pemrograman Internet yang berbasis Outcome-Based Education (OBE). Keterkaitan ini diwujudkan melalui pemetaan eksplisit setiap bab ke Sub-CPMK dan CPMK yang telah ditetapkan dalam silabus. Pendekatan ini memastikan bahwa materi pembelajaran, aktivitas, dan asesmen sejalan dengan capaian yang diharapkan oleh Program Studi Teknik Informatika, Fakultas Teknologi Informasi, Universitas Bale Bandung.

\textbf{Alignment dengan CPL dan CPMK:} Setiap bab dalam buku ini dipetakan ke Sub-CPMK yang berkontribusi pada pencapaian CPMK dan CPL. Bab II--IV mencakup Sub-CPMK 1.1, 1.2 (arsitektur web) dan Sub-CPMK 2.1 (HTML5). Bab V--VIII mencakup Sub-CPMK 2.2, 2.3 (CSS3 dan JavaScript). Bab IX--XIV mencakup Sub-CPMK 3.1, 3.2 (server-side, database) dan Sub-CPMK 4.1, 4.2 (validasi input, session, keamanan).

\textbf{Integrasi Metode Pembelajaran:} Buku ini mengintegrasikan metode pembelajaran yang tercantum dalam RPS, meliputi ceramah interaktif (30\% teori), praktikum terbimbing (70\% praktek), demonstrasi dan live coding, Problem-Based Learning (PBL), Project-Based Learning, peer review, dan studi kasus. Setiap bab menyediakan aktivitas yang mendukung metode tersebut sesuai konteks materi \cite{freecodecamp}.

\textbf{Sistem Penilaian:} Komponen asesmen dalam buku ini sejalan dengan bobot penilaian RPS: Tugas Individu (15\%), Kuis (10\%), Praktikum (20\%), UTS (20\%), Proyek Kelompok (20\%), dan UAS (15\%).

\begin{table}[h]
\centering
\small
\begin{tabular}{|c|p{5.5cm}|c|}
\hline
\textbf{Bab} & \textbf{Topik} & \textbf{Sub-CPMK} \\
\hline
II & Landasan Teori & 1.1, 1.2 \\
\hline
III--IV & HTML5 & 2.1 \\
\hline
V--VI & CSS3 & 2.2 \\
\hline
VII--VIII & JavaScript & 2.3 \\
\hline
IX--X & PHP, MySQL & 3.1, 3.2 \\
\hline
XI--XIV & Session, integrasi, keamanan, evaluasi & 4.1, 4.2 \\
\hline
\end{tabular}
\caption{Pemetaan Bab ke Sub-CPMK}
\end{table}
